% !TEX root = ../main.tex
\CourseSection{Design and Production Approach}\label{sec:design}

\CourseSubsection{Metric Development and Analysis}

The following outlines the proposed methods for measuring size, churn, age, clone coverage, and readability, including the data requirements, calculation procedures, and validation approach for each metric.
Each metric follows four rough stages: it is prototyped on the official 13,000-skill GitSkills sample, calibrated, scaled to the full dataset, and validated.

\noindent\textbf{Size.}
Size is classically measured in lines~\cite[ch.~8]{fenton2014metrics}, and the closest prior work counted words in context files~\cite{chatlatanagulchai2026agent}, but neither captures what an agent pays for.
A skill mixes plain-language instructions with code blocks, and a line of prose and a line of code carry very different amounts of information; in our pilot, about 12\% of a typical skill's body was code.
Instead, we measure size in tokens, the units an agent actually reads text in and the units its context window is limited by, counted with Anthropic's tokenizer.
We split that count across the three parts of a skill, since an agent loads each at a different time: the \textit{description}, read for every installed skill to decide which one to use; the \textit{body}, loaded only once a skill is chosen; and the \textit{bundle}, the scripts and reference files stored next to the skill, which the body points to when it needs them.
Measuring each part separately shows not just how big a skill is, but where that size costs the agent.
Lines and words are still recorded alongside tokens, so skills can be compared directly with the code bundled next to them.\par

\noindent\textbf{Churn.}
Classically, churn is the number of lines added plus lines removed between commits, and relative churn has been shown to predict defect density~\cite{nagappan2005churn}.
Since lines don't fit skills well, churn needs adapting too, and the data adds a second complication: GitSkills records how often and when a skill was committed, but not what changed, and about two-thirds of skills in our pilot were never edited after being added.
What skills tend to do instead is get copied and then modified for a new project, so we split churn into three types.
\textit{Activity churn} uses the recorded commit counts and dates to measure how often, and for how long, a skill is maintained.
\textit{Cross-copy churn} takes the near-clone families found by Clone Coverage (MinHash), links each skill to the most similar member that appeared before it, and counts the tokens added and removed along the way. Which copy came first is decided by the first commit date, or by the repository's creation date when a skill has no history, and a sample of the links is checked by hand (PF-04). This is where we expect the most useful results, since it captures how popular templates get adapted for uses like design or security. A copy creates an implicit dependency rather than a declared one: nothing records the link, and the child receives no later changes from its parent, the situation clone genealogies describe~\cite{kim2005genealogies}. We therefore report links as lineage and do not claim that a child relies on its parent at run time.
\textit{Token churn} is classical churn counted in tokens instead of lines, computed from the full git history of a sample of edited skills, and confirms that the other two measures reflect real change.\par

\noindent\textbf{Age.}
Age measures how long a skill has existed within its repository, following work that uses module age and change history to predict faults~\cite{graves2000predicting}. The GitSkills \texttt{artifacts} table provides the \texttt{first\_commit\_at} and \texttt{last\_commit\_at} timestamps, of which \texttt{first\_commit\_at} will be used to calculate age. For each eligible skill, we will subtract its first recorded commit date from a fixed dataset collection date and express the elapsed time in days. Using the same collection date for every skill ensures that the results are reproducible across repeated runs.

The first commit is preferred over the last commit because the project measures how long a skill has existed, not how long it was actively maintained. For example, a skill that has never been edited after its introduction would still have a meaningful age. Since skills may be copied between repositories, the first recorded commit represents the skill's observed introduction to its current repository rather than necessarily the original creation date of its instructions. Records with missing, invalid, or future-dated timestamps will be flagged rather than assigned an estimated age. The calculation will first be checked against manually calculated examples on the 13,000-skill sample, then applied to eligible records in the full dataset.\par

\noindent\textbf{Clone Coverage.}

Clone coverage measures the proportion of skill occurrences belonging to groups with identical or substantially similar content, adapting the code clone literature~\cite{roy2009comparison,kamiya2002ccfinder}. We report it at file level rather than line or token level because skills spread by copying whole folders, because \texttt{file\_sha} makes file-level duplication exact and reproducible, and because a file-level map has already described duplication across GitHub~\cite{lopes2017dejavu}. Line-level coverage would also split a skill's instructions into fragments that an agent never loads alone. Near-clones, measured on token shingles~\cite{broder1997resemblance}, capture the partial overlap that file hashes miss. Exact clones will be identified in the GitSkills \texttt{artifacts} table by grouping occurrences with matching \texttt{file\_sha} values. Any hash group containing more than one occurrence will be classified as an exact clone group. Exact clone coverage will be calculated as the number of occurrences in these groups divided by the total number of eligible occurrences, multiplied by 100. Every occurrence in a duplicated group is counted, including the first one. The \texttt{repo\_full\_name} and \texttt{path} fields identify where occurrences appear, but will not be used to determine text similarity.

Near-clones will also be detected using the actual skill text in \texttt{content}. The \texttt{dedup\_primary} field will help select one representative per exact-content group to avoid repeated comparisons, while \texttt{content\_fetched} and \texttt{content\_sha\_ok} will be checked for text availability and integrity. Each representative's text will be split into overlapping three-token sequences (shingles). MinHash signatures and locality-sensitive hashing (LSH) will retrieve candidate pairs with potentially similar content; their actual Jaccard similarity will then be computed as the number of shared shingles divided by the number of distinct shingles across both skills. Pairs with different file hashes and Jaccard similarity meeting a selected threshold will be linked as near-clones. Rather than fixing the threshold beforehand, several candidate thresholds will be evaluated using manually reviewed skill pairs to determine which best distinguishes substantial text reuse from unrelated skills. This is a text-reuse criterion, not proof that one repository copied from another.

Exact clone groups and qualifying near-clone links will be combined into connected groups. This permits chains of related skills even when not every pair in the group individually meets the selected similarity threshold, so the grouping rule will be reported explicitly. Near-clone relationships will be mapped from representative contents back to all eligible skill occurrences using \texttt{file\_sha}. We will report exact clone coverage and expanded coverage including near-clones separately, with the latter calculated over occurrences for which usable similarity data is available. Missing hashes, unavailable content, and exclusions will be reported so that the coverage percentages can be interpreted correctly. The methods will be prototyped on the 13,000-skill sample before scaling to the full dataset.\par


\noindent\textbf{Readability.}
Section~\ref{sec:functional-description} describes the module, which we prototyped on the 13,000-skill sample: median skill bodies read at grade 9.1, descriptions at grade 12.7, and about 19\% of eligible skills are not in English (Appendix~\ref{app:readability-pilot}). The next iteration replaces the pilot's simple English check with fastText~\cite{joulin2017bag}, followed by the hand review (T05.5).\par

\CourseSubsection{Milestones and Division of Labor}

Development will follow an iterative approach: metric implementations will be tested on sample data and refined before integration and larger-scale analysis. Progress and outstanding issues will be reviewed during Friday team meetings. Each milestone lead will coordinate completion and supporting test evidence, with implementation and review shared across the group.

Across those iterations, the prototypes in Table~\ref{tab:prototypes} add capability in increments: one stage, then one module at a time, then the integrated pipeline. Modules join P3 only after their T01 to T05 tests pass on fixtures, and the join key (\texttt{file\_sha}) is checked by T06 at each integration.

\begin{singlespace}
\setlength{\LTcapwidth}{\linewidth}
\begin{longtable}{@{}L{0.07\linewidth} L{\dimexpr0.47\linewidth-6\tabcolsep\relax} L{0.25\linewidth} L{0.21\linewidth}@{}}
\caption{Prototypes, with success criteria and traceability.}\label{tab:prototypes}\\
\rowcolor{MidGray}
\TableHeader{ID} & \TableHeader{Purpose, scope, and development} & \TableHeader{Success criteria} & \TableHeader{Tests; milestones}\\
\midrule
\endfirsthead
\rowcolor{MidGray}
\TableHeader{ID} & \TableHeader{Purpose, scope, and development} & \TableHeader{Success criteria} & \TableHeader{Tests; milestones}\\
\midrule
\endhead
\bottomrule
\endfoot
P1 & Show that the data loads and filters correctly. Covers F1 and F2 on the 13,000-skill sample, starting from the pilot filters. & IF-01, IF-02, SYS-2 & T06; M2\\
\addlinespace[3pt]
P2 & Show that each metric calculates correctly alone. Covers F3.1 to F3.5, on fixtures and then the sample, one module at a time. & SYS-1, PF-04, and each metric's requirements & T01 to T05; M1, M3\\
\addlinespace[3pt]
P3 & Show that the modules work together and scale. Covers F4 and F5, first on the sample and then on the full dataset in batches. & IF-03, PF-01 to PF-03, SYS-3, SYS-5, SYS-6 & T06 to T10; M4, M5, M6, M12\\
\end{longtable}
\end{singlespace}

\begin{singlespace}
\setlength{\LTcapwidth}{\linewidth}
\begin{longtable}{@{}L{0.09\linewidth} L{\dimexpr0.55\linewidth-6\tabcolsep\relax} L{0.18\linewidth} L{0.18\linewidth}@{}}
\caption{Planned milestones and lead responsibilities.}\label{tab:proposal-milestones}\\
\rowcolor{MidGray}
\TableHeader{ID} & \TableHeader{Milestone} & \TableHeader{Target date} & \TableHeader{Lead}\\
\midrule
\endfirsthead
\rowcolor{MidGray}
\TableHeader{ID} & \TableHeader{Milestone} & \TableHeader{Target date} & \TableHeader{Lead}\\
\midrule
\endhead
\bottomrule
\endfoot

M1 & Metric definitions reviewed
& Nov.\ 6, 2026 & Noah\\
\addlinespace[3pt]

M2 & Dataset preparation completed
& Nov.\ 13, 2026 & Naman\\
\addlinespace[3pt]

M3 & Initial metric calculations verified
& Nov.\ 13, 2026 & James\\
\addlinespace[3pt]

M4 & Sample pipeline integrated
& Nov.\ 20, 2026 & Faiz\\
\addlinespace[3pt]

M5 & Scalability assessed
& Nov.\ 27, 2026 & Naman\\
\addlinespace[3pt]

M6 & Initial analysis produced
& Dec.\ 4, 2026 & Faiz\\
\addlinespace[3pt]

M7 & MSR submission-readiness review
& Dec.\ 11, 2026 & James\\
\addlinespace[3pt]

M8 & MSR abstract submitted (conditional)
& Dec.\ 18, 2026 & Faiz\\
\addlinespace[3pt]

M9 & MSR paper submitted (conditional)
& Dec.\ 23, 2026 & Noah\\
\addlinespace[3pt]

M10 & Results reviewed and focused research questions documented
& Jan.\ 29, 2027 & Noah\\
\addlinespace[3pt]

M11 & Follow-up investigations evaluated
& Feb.\ 26, 2027 & Naman\\
\addlinespace[3pt]

M12 & Final pipeline validated and results reproduced
& Mar.\ 12, 2027 & James\\

\end{longtable}
\end{singlespace}

M2 and M3 are in progress. Dataset preparation has supported preliminary assessments, but final preprocessing checks and initial calculation verification remain ongoing.

M1 establishes the definitions used for M3. Completion of M2 and M3 supports integration in M4, followed by scalability assessment in M5 and initial analysis in M6. M10 reviews the results and selects follow-up investigations for M11 before final validation in M12.

M7 will assess findings, validation evidence, limitations, and writing readiness to determine whether to proceed with conference submission. M8 proceeds if that review supports submission; M9 additionally depends on sufficiently validated findings and supporting analysis. The abstract and paper deadlines are December 18 and December 23, 2026, respectively~\cite{msr2027MiningChallenge}. M8 and M9 are conditional; the remaining capstone milestones apply regardless of the submission decision.

\CourseSubsection{Further Exploration and Contingencies}\label{sec:contingencies}

Following the initial analysis in M6, the research review in M10 will select follow-up investigations for M11. Candidates include differences in skill revision activity across repository domains, instruction structure, reuse patterns, and an LLM-assisted taxonomy of skill purposes~\cite{wan2024tnt}. Selection will depend on data availability, measurement validity, and remaining time.

\begin{singlespace}
\setlength{\LTcapwidth}{\linewidth}
\hyphenpenalty=10000 \exhyphenpenalty=10000
\begin{longtable}{@{}L{0.05\linewidth} L{\dimexpr0.27\linewidth-2\tabcolsep\relax} L{\dimexpr0.52\linewidth-2\tabcolsep\relax} L{\dimexpr0.16\linewidth-2\tabcolsep\relax}@{}}
\caption{Contingencies and affected milestones.}\label{tab:contingencies}\\
\rowcolor{MidGray}
\TableHeader{ID} & \TableHeader{Risk and trigger} & \TableHeader{Response and effect on the plan} & \TableHeader{Affected milestones}\\
\midrule
\endfirsthead
\rowcolor{MidGray}
\TableHeader{ID} & \TableHeader{Risk and trigger} & \TableHeader{Response and effect on the plan} & \TableHeader{Affected milestones}\\
\midrule
\endhead
\bottomrule
\endfoot

R1 & \textbf{Incomplete data or repository history:} fields or revisions required for a calculation cannot be obtained.
& Document missingness and restrict the affected analysis to eligible records. If insufficient evidence remains, report the measurement as unavailable and assess the effect on requirements and conclusions.
& M2, M3, M6\\
\addlinespace[3pt]

R2 & \textbf{Unsupported metric interpretation:} a metric fails its specified validation criteria.
& Review preprocessing, assumptions, and the measurement method, then repeat validation. If it remains unsupported, report its limitations and exclude interpretations that the evidence does not justify.
& M3, M6, M10, M12\\
\addlinespace[3pt]

R3 & \textbf{Processing limits:} measured runtime or memory exceeds the specified acceptance limit (PF-01, PF-02).
& Adjust batch sizes or data access and repeat performance tests. If full-dataset processing remains infeasible, propose documented sampling and explicitly review its effect on coverage and performance requirements.
& M5, M6, M12\\
\addlinespace[3pt]

R4 & \textbf{Schedule pressure:} unresolved defects or analysis delays threaten required validation or delivery.
& Reassign work and defer optional investigations to protect core analysis and validation. Forgo MSR submission if necessary and continue toward capstone deadlines.
& M6, M7, M8, M9, M11, M12\\
\addlinespace[3pt]

R5 & \textbf{Cost:} storage or compute may exceed team hardware, which would break the \$0 budget. Trigger: PF-01 or PF-02 is missed after R3.
& Request access to Queen's computing resources, as our supervisor suggested, before considering any paid service. Budget stays at \$0 unless the team and supervisor agree otherwise.
& M5, M6\\
\addlinespace[3pt]

R6 & \textbf{Programmatic:} the dataset release, GitHub access, or challenge rules may change outside our control. Trigger: a new release, or clones fail for token churn.
& Keep a local copy of the July 2026 snapshot and record its version. Report token churn on the histories that can be retrieved, and review any rule change at M7.
& M2, M3, M7\\

\end{longtable}
\end{singlespace}

Changes affecting mandatory requirements will be reviewed with the supervisor and reflected in the requirements, tests, and schedule. Optional investigations may be deferred, but unsuccessful tests or unsupported findings will not be presented as completed validation.
